\input{preamble}

\title{Section 1.4 --- Inverse Relations --- Answer Key}

\begin{document}
\maketitle

\section{Writing Inverses}
Write the inverse of each relation.
\begin{smenum}
\mitemx
{$P\inv=\Set{(-6,4),(-2,-10),(-9,-1),(-7,0),(-9,3)}$}
\mitemx
{$Q\inv=\Set{(5,-7),(3,-2),(9,-2),(-5,10)}$}
\mitemx
{$R\inv=\Set{(11,-2),(-8,-1),(3,-9),(5,6),(-2,-8)}$}
\mitemx
{$S\inv=\Set{(-10,-10),(1,5),(-3,3),(9,9),(-6,3)}$}
\mitemxx
{$T\inv=\Relat{x=3y-4}$}
{$U\inv=\Relat{y^2=x^3}$}
\mitemxx
{$V\inv=\Relat{y=x^5+y^2x+x^3}$}
{$W\inv=\Relat{y^2+x^2=25}$}
\end{smenum}

\section{Graphing Inverses}
Several relations are graphed below.  Sketch the graphs of their inverses.  \emph{Hint: Used the points marked to guide your sketch.}

(Note: the original graph is shown grayed out for comparison.)
\begin{smenum}
\mitemxx
{\LittleGraphEO[.15in]{
	\draw[thick,domain=-4:4,smooth,black!25!white] plot (1/3*\x*\x*\x-4/3*\x,-2*\x);
	\foreach \x in {-2.5,-2,-1,0,1,2,2.5}
			\fill[black!25!white] (1/3*\x*\x*\x-4/3*\x,-2*\x) circle (.2);
	\draw[thick,domain=-4:4,smooth] plot (-2*\x,1/3*\x*\x*\x-4/3*\x);
	\foreach \x in {-2.5,-2,-1,0,1,2,2.5}
			\fill (-2*\x,1/3*\x*\x*\x-4/3*\x) circle (.2);
}}
{\LittleGraphEO[.15in]{
	\draw[thick,domain=-5.9:5.9,black!25!white] plot (\x,.5*\x*\x-3);
	\foreach \x in {-4,-2,...,4}
		\fill[black!25!white] (\x,.5*\x*\x-3) circle (.2);
	\draw[thick,domain=-5.9:5.9] plot (.5*\x*\x-3,\x);
	\foreach \x in {-4,-2,...,4}
		\fill (.5*\x*\x-3,\x) circle (.2);
}}
\mitemxx
{\LittleGraphEO[.15in]{
    \draw[thick,black!25!white] (-5.9,3/4*5.9+2) -- (5.9,-3/4*5.9+2);
		\foreach \x in {-4,0,4}
		    \fill[black!25!white] (\x,-3/4*\x+2) circle (.2);
    \draw[thick] (3/4*5.9+2,-5.9) -- (-3/4*5.9+2,5.9);
		\foreach \x in {-4,0,4}
		    \fill (-3/4*\x+2,\x) circle (.2);
}}
{\LittleGraphEO[.15in]{
	\draw[thick,xscale=.75,black!25!white] (4/3,1) circle (4);
	\foreach \x/\y in {1/5,1/-3,4/1,-2/1}
		\fill[black!25!white] (\x,\y) circle (.2);
	\draw[thick,yscale=.75] (1,4/3) circle (4);
	\foreach \x/\y in {1/5,1/-3,4/1,-2/1}
		\fill (\y,\x) circle (.2);
}}
\end{smenum}
\section{Domain and Range of the Inverse}
\begin{senum}
\item $\dom R=\Set{-2,3}$, $\rng R=\Set{4,5,6}$, $\dom R\inv=\Set{4,5,6}$, $\rng R\inv=\Set{-2,3}$
\item $\dom S\inv=\rng S=\Rint{11}$, $\rng S\inv=\dom S=\Lrint{-5}{2}$
\item $\rng T=\dom T\inv=\SetBuild{x\neq -1}$, $\rng T\inv=\dom T=\SetBuild{x\neq 5}$
\end{senum}

\clearpage
\section{Identifying One-to-One}
For each relation below, decide which of the following applies to the relation. (This is multiple choice.)
\begin{clist}
\item A function (but not one-to-one)
\item One-to-one but not a function
\item A one-to-one function
\item Not one-to one or a function
\end{clist}
\begin{smenum}
\mitemxx
{a - a function (but not one-to-one)}
{d - not one-to-one or a function}
\mitemxx
{a - a function (but not one-to-one)}
{b - one-to-one but not a function}
\mitemxx
{d - not one-to-one or a function}
{a - a function (but not one-to-one)}
\mitemxx
{c - a one-to-one function}
{b - one-to-one but not a function}
\end{smenum}
For each relation below, determine which of the four ordered pairs given are in the relation. Use this to decide what kind of relation this is. (This is multiple choice.)
\begin{clist}
\item A function (but not one-to-one)
\item One-to-one but not a function
\item A one-to-one function
\item Not one-to one or a function
\end{clist}
The pairs are chosen so that, if the relation is \emph{not} one-to-one, the given ordered pairs will demonstrate this.
\begin{smenum}
\mitemxx
{c - a one-to-one function}
{d - not one-to-one or a function}
\mitemxx
{a - a function (but not one-to-one)}
{b - one-to-one but not a function}
\end{smenum}
\section{Inverse Functions}
The function $h=\Set{(-4,-7),(2,2),(4,8),(8,-9)}$ is one-to-one. Find each of the following, or write ``undefined.''
\begin{smenum}
\mitemxxxx
{$h\inv(-7)=-4$}
{$h\inv(8)$ is undefined}
{$h\inv(4)$ is undefined}
{$h\inv(2)=2$}
\end{smenum}
For each function given below, write the inverse function. Answer in function notation.
\begin{smenum}
\mitemxx
{$f\inv(x)=\frac{1}{2}x+4$}
{$g\inv(x)=\frac{1}{5}x-3$}
\end{smenum}
The function $k(x)=\dfrac{x-3}{2x+1}$ is one-to-one. Evaluate each of the following.
\begin{smenum}
\mitemxxxx
{$k(-4)=1$}
{$k(-1)=4$}
{$k(0)=-3$}
{$k(3)=0$}
\end{smenum}
Based on the values of $k(x)$ you found, evaluate each of the following.
\begin{smenum}
\mitemxxxx
{$k\inv(-3)=0$}
{$k\inv(0)=3$}
{$k\inv(1)=-4$}
{$k\inv(4)=-1$}
\end{smenum}

\end{document}